\originalchapter{计数、样本空间与概率公理}
\sourcelecture{1--5}

概率论的第一项工作不是代入公式, 而是说明一次试验究竟记录什么. 同样是抽取若干物品, 记录抽取顺序与只记录最后的集合, 得到的是不同的样本空间. 若两种记录之间每个结果恰好有同样多的对应关系, 两种算法可以给出相同的概率; 若对应数不同, 就不能把粗略记录后的结果一律看成等可能. 本章先建立计数工具, 再说明什么样的数字分配可以称为概率.

\originalsection{分阶段选择与排列}
设完成一件事需要依次作出若干选择. 若第一步有 $a_1$ 种选择, 而无论前面如何选择, 第 $j$ 步总有 $a_j$ 种选择, 则总数为 $a_1\cdots a_m$. 这是乘法原理. 它可由一棵选择树解释: 每个第一层结点有同样多的后继, 逐层数到叶子即可. 条件中的``无论前面如何选择''不能省略. 如果某一步的选择数会变化, 应先分类, 对各类求乘积, 再把互不重叠的各类相加.

\begin{example}
用符号 $0,1,2,3$ 写长度为六的序列, 要求相邻符号不同. 若另要求首位不能为零, 各有多少个?
\end{example}
\begin{solution}
第一种情形首位有四种选择. 后面的每一位只须避开紧邻的前一位, 因而总有三种选择, 答案为 $4\cdot3^5$. 第二种情形只有首位的选择数变为三, 答案为 $3^6$. 后一位究竟可用哪些符号会随前一位变化, 但可用符号的\emph{数量}始终为三; 乘法原理要求的是数量恒定.
\end{solution}

从 $n$ 个不同元素中有序选出 $k$ 个, 不允许重复, 总数为
\[
 n(n-1)\cdots(n-k+1)=\frac{n!}{(n-k)!},\qquad 0\leq k\leq n.
\]
允许重复时则为 $n^k$. 这里 $n!=1\cdot2\cdots n$, 并约定 $0!=1$. 空序列只有一个, 从空集合到自身也只有一个函数, 所以这个约定表达了真实的计数事实, 同时使边界 $k=0,n$ 无须另改公式.

\begin{definition}
集合 $\{1,\ldots,n\}$ 的排列是该集合到自身的双射 $\sigma$. 用序列 $\sigma(1),\ldots,\sigma(n)$ 可以表示它, 因而排列总数为 $n!$. 两个排列的复合定义为 $(\sigma\circ\rho)(j)=\sigma(\rho(j))$.
\end{definition}
排列不只是把元素排成一列的记号. 反复施用 $\sigma$ 可以看出它的结构. 从元素 $j$ 出发, 依次写 $j,\sigma(j),\sigma^2(j),\ldots$. 有限性保证某时重复; 双射性保证第一次重复必定回到 $j$, 否则对重复位置施用逆映射会产生更早的重复. 得到的闭合序列叫一个循环. 对尚未出现的元素重复操作, 就把集合分成互不相交的循环. 循环内部换一个起点、交换各循环书写次序, 不改变排列.

例如 $(1\ 4\ 2)(3\ 5)(6)$ 表示 $1\mapsto4\mapsto2\mapsto1$, $3\leftrightarrow5$, 而 $6$ 不变. 满足 $\sigma^2(j)=j$ 的排列叫对合. 它恰好是循环长度只有一或二的排列: 循环长度为 $r$ 时, 施用两次回到原点等价于 $r$ 整除二. 长度为一的循环称固定点. 没有固定点的排列称错排, 我们将在容斥原理之后计算其数量.

\originalsection{消除顺序与组合数}
从 $n$ 个不同元素中选出 $k$ 个组成集合, 与有序选择的区别在于同一集合的 $k!$ 种次序全部合并. 每个集合恰好对应 $k!$ 个有序序列, 因而
\[
 \binom nk=\frac{n!}{k!(n-k)!}.
\]
这种``先数带标签的对象, 再除以固定重复倍数''的方法很常用. 除法的依据必须是每个最终对象都有相同的原像数. 若把两个颜色相同的物品看成不可区分, 就应除以它们内部的排列数; 若不同对象有不同数目的对称性, 则不能直接除以一个猜出的常数.

\begin{proposition}[Pascal 递推与二项式定理]
约定 $\binom nk=0$ 当 $k<0$ 或 $k>n$. 对 $n\geq1$,
\[
 \binom nk=\binom{n-1}{k-1}+\binom{n-1}{k},\qquad
 (x+y)^n=\sum_{k=0}^n\binom nk x^ky^{n-k}.
\]
特别地, $\sum_{k=0}^n\binom nk=2^n$.
\end{proposition}
\begin{proof}
固定一个指定元素. 含有它的 $k$ 元子集, 对应于从其余 $n-1$ 个元素选 $k-1$ 个; 不含它的, 对应于选 $k$ 个. 两类互不重叠且穷尽所有选择, 得递推式.

展开 $n$ 个因子 $(x+y)$ 的乘积时, 每一项来自在每个因子中选 $x$ 或 $y$. 要得到 $x^ky^{n-k}$, 恰好要指定哪 $k$ 个因子选 $x$, 所以系数为 $\binom nk$. 令 $x=y=1$ 得最后的等式. 这也说明 $n$ 元集合共有 $2^n$ 个子集, 因为每个元素独立地有``选入''和``不选入''两种计数选择.
\end{proof}

\begin{example}
十张卡片具有两种类别, 每类各五张且每张有不同编号. 不计顺序地选四张, 要求恰有三张来自第一类. 若改成逐张抽取并记录顺序, 如何验证同一比例?
\end{example}
\begin{solution}
无序选择的总数为 $\binom{10}{4}$, 合格数为 $\binom53\binom51$. 有序总数为 $10\cdot9\cdot8\cdot7$. 合格有序数为
\[
 \binom43(5\cdot4\cdot3)\cdot5.
\]
这里先指定三张第一类卡片出现的位置, 再依次填入卡片. 前后两种总数及合格数分别相差 $4!$, 因而比例相同. 不能把有序合格数除以无序总数.
\end{solution}

\originalsection{多项式系数与整数分配}
把 $n$ 个不同元素分给 $r$ 个有名称的组, 组大小依次是非负整数 $n_1,\ldots,n_r$, 且 $n_1+\cdots+n_r=n$. 把全部元素先排列, 再按规定大小切成各组, 同一分组被各组内部排列重复 $n_1!\cdots n_r!$ 次. 所以分组数为
\[
 \binom{n}{n_1,\ldots,n_r}=\frac{n!}{n_1!\cdots n_r!}.
\]
组的名称很重要: 给甲、乙两个部门分配与把人分成两个无名称的团队是不同问题. 当无名称组的大小相同, 还会有交换这些组产生的重复.

\begin{theorem}[多项式定理]
对非负整数 $n$,
\[
 (x_1+\cdots+x_r)^n
 =\sum_{\substack{n_1,\ldots,n_r\geq0\\n_1+\cdots+n_r=n}}
 \frac{n!}{n_1!\cdots n_r!}x_1^{n_1}\cdots x_r^{n_r}.
\]
\end{theorem}
\begin{proof}
乘积的每个展开项对应一个长度为 $n$ 的符号序列, 每个位置选 $r$ 个符号之一. 若符号 $i$ 出现 $n_i$ 次, 消去位置标签之后这一项变为 $x_1^{n_1}\cdots x_r^{n_r}$. 同一单项式出现的次数, 正是把 $n$ 个位置分成各个符号所占位置组的数量, 即上面的多项式系数. 对所有可能的次数向量求和即可.
\end{proof}
令所有 $x_i=1$, 可知多项式系数之和为 $r^n$. 比如 $x^2yz^3$ 在 $(x+y+z)^6$ 中的系数为 $6!/(2!1!3!)=60$.

\begin{proposition}[隔板计数]
对整数 $n\geq0,k\geq1$, 非负整数解
\[
 a_1+\cdots+a_k=n
\]
共有 $\binom{n+k-1}{k-1}$ 个. 若各 $a_i\geq1$ 且 $n\geq k$, 则共有 $\binom{n-1}{k-1}$ 个.
\end{proposition}
\begin{proof}
写出 $n$ 个星号和 $k-1$ 条隔板. 第一条隔板之前的星号数为 $a_1$, 相邻隔板之间的数依次为 $a_2,\ldots,a_{k-1}$, 最后为 $a_k$. 隔板可以相邻, 也可位于两端, 正好允许零值. 这个对应可以反过来唯一恢复星号隔板串. 从 $n+k-1$ 个位置中选出隔板位置, 得第一式. 正整数情形令 $b_i=a_i-1$, 就变为非负整数和 $n-k$ 的问题, 因而数量为 $\binom{n-1}{k-1}$.
\end{proof}
这里数的是有序的大小向量, 又叫弱组合, 不是通常将不同次序视为同一对象的整数分拆. 也不要把它与多项式系数混淆: 隔板数的是可能的组大小; 多项式系数是在组大小已给定时分配不同元素的数量.

\begin{remark}[阶乘与积分]
微积分还给出一个与以后连续分布相连的表达式. 设 $I_n=\int_0^\infty t^n\ee^{-t}\dd t$. 当 $n=0$ 时 $I_0=1$. 对 $n\geq1$ 分部积分, 边界项 $-t^n\ee^{-t}$ 在零和无穷处都为零, 所以 $I_n=nI_{n-1}=n!$. 对实数 $z>0$, 积分 $\Gamma(z)=\int_0^\infty t^{z-1}\ee^{-t}\dd t$ 收敛, 且同一计算给 $\Gamma(z+1)=z\Gamma(z)$. 因而 $\Gamma(n+1)=n!$. 这是一种把阶乘延伸到非整数的方式, 本章的计数仍只使用整数阶乘.
\end{remark}

\originalsection{样本空间与事件}
一次试验所有可能记录组成的集合称为样本空间, 记为 $\Omega$. 其中一个元素 $\omega$ 是一个完整结果. 事件是我们能够询问``是否发生''的结果集合. 掷两枚有标号的骰子时, 可用 $\Omega=\{1,\ldots,6\}^2$; ``点数和为四''是集合 $\{(1,3),(2,2),(3,1)\}$. 若只记录点数和, 空间为 $\{2,\ldots,12\}$, 但这些和并不等可能.

事件 $A\cup B$ 表示至少一个发生, $A\cap B$ 表示两者都发生, $A^c=\Omega\setminus A$ 表示 $A$ 不发生. 有时简写 $AB=A\cap B$. $A\setminus B=A\cap B^c$. 两事件互斥是说 $A\cap B=\varnothing$, 即不可能同时发生. 对任意事件族, De Morgan 法则为
\[
 \left(\bigcup_i A_i\right)^c=\bigcap_i A_i^c,\qquad
 \left(\bigcap_i A_i\right)^c=\bigcup_i A_i^c.
\]
证明只须逐个结果判断: 不在任何 $A_i$ 中, 等价于在每个 $A_i^c$ 中; 不在全部 $A_i$ 中, 等价于至少在一个补集中. 因而这些等式既适用于有限族, 也适用于可数族.

在有限空间上, 所有子集都可以作为事件. 在连续空间上, 若要求概率满足合理的可数可加性, 通常不能对任意子集都赋予一致的长度或面积概率. 我们因此明确指定一个允许的事件族.

\begin{definition}
$\Omega$ 上的 $\sigma$ 代数是子集族 $\mathcal F$, 满足 $\Omega\in\mathcal F$, 若 $A\in\mathcal F$ 则 $A^c\in\mathcal F$, 若 $A_1,A_2,\ldots\in\mathcal F$ 则 $\bigcup_{n=1}^\infty A_n\in\mathcal F$. 概率空间是三元组 $(\Omega,\mathcal F,\Prob)$, 其中 $\Prob:\mathcal F\to[0,1]$ 满足
\[
 \Prob(\Omega)=1,\qquad
 \Prob\left(\bigcup_{n=1}^\infty A_n\right)=\sum_{n=1}^\infty\Prob(A_n)
 \quad\text{当 }A_n\text{ 两两互斥}.
\]
\end{definition}
补集与并集的封闭性, 通过 De Morgan 法则也给出对可数交集的封闭性. 本书有限、可数模型通常取 $\mathcal F$ 为全部子集. 对实数上的连续模型, 使用包含所有区间、并对上述操作封闭的 Borel 事件族, 必要时再补入零概率集的子集. 本书不证明不可测集合的存在, 也不要求读者构造它们; 这里说明事件必须属于 $\mathcal F$, 防止把有限模型的方便约定误当成对所有连续子集都成立的结论.

概率公理给出数学的一致性, 并不单独决定现实中各事件的概率. 长期频率、主观信念和市场价格可以帮助建立模型, 但各有前提. 重复试验的经验比例满足有限集合运算的一致性, 而把未来概率等同于过去频率还需稳定性假设. 若以合约价格解释概率, 需要同一计价单位及可以交易相应合约等条件. 有限次无套利比较解释有限可加性, 本身不推出可数可加性. 个人认为某个故事更可信, 也不等于它的事件集合更大.

\begin{example}
某报价系统给互斥事件 $A,B$ 的一单位支付合约分别定价 $0.2,0.3$, 却给 $A\cup B$ 的一单位支付合约定价 $0.7$. 假定可以买入和卖出这些合约, 且不计手续费. 说明这些价格为什么不一致.
\end{example}
\begin{solution}
买入 $A$ 和 $B$ 两份合约的总成本为 $0.5$, 卖出 $A\cup B$ 合约收到 $0.7$, 当场净收入为 $0.2$. 因 $A,B$ 互斥, 两份买入合约的总支付在每个结果上都恰好等于卖出合约的支付, 未来净支付为零. 因而留下确定的 $0.2$ 收入. 这解释了在允许双向交易的理想化系统中, 有限可加性如何来自相同支付必须有相同价格. 若合约交易有限制或交易成本不可忽略, 这种推理的假设就须重新检查.
\end{solution}

\originalsection{从公理导出运算法则}
\begin{proposition}
概率空间中有 $\Prob(\varnothing)=0$, $\Prob(A^c)=1-\Prob(A)$. 若 $A\subseteq B$, 则
\[
 \Prob(B)=\Prob(A)+\Prob(B\setminus A)\geq\Prob(A).
\]
此外,
\[
 \Prob(A\cup B)=\Prob(A)+\Prob(B)-\Prob(A\cap B).
\]
\end{proposition}
\begin{proof}
将 $\Omega$ 与无穷多个空集作为互斥事件代入可数可加性, 得 $1=1+\sum_{n\geq1}\Prob(\varnothing)$, 非负性迫使空集概率为零. 因而可数可加性包含有限可加性. $\Omega=A\cup A^c$ 是互斥分解, 得补集公式. 对 $A\subseteq B$, 使用 $B=A\cup(B\setminus A)$ 的互斥分解, 得单调性. 最后 $A\cup B=A\cup(B\setminus A)$, 而 $B=(B\setminus A)\cup(A\cap B)$, 两式相减即得并集公式.
\end{proof}
单调性意味着 $\Prob(A\cap B)\leq\Prob(A)$. 例如``某人从事某职业并且参加某社团''不可能比``某人从事某职业''概率更大, 无论附加细节怎样使故事听起来更贴合描述. 这是事件包含关系的后果.

\begin{proposition}[连续性与并集上界]
若 $A_n\subseteq A_{n+1}$, 则 $\Prob(\bigcup_n A_n)=\lim_n\Prob(A_n)$. 若 $A_{n+1}\subseteq A_n$, 则 $\Prob(\bigcap_n A_n)=\lim_n\Prob(A_n)$. 对任意可数事件族,
\[
 \Prob\left(\bigcup_n A_n\right)\leq\sum_n\Prob(A_n).
\]
\end{proposition}
\begin{proof}
递增情形令 $D_1=A_1$, $D_n=A_n\setminus A_{n-1}$, 则 $D_n$ 两两互斥, $A_n=\bigcup_{j\leq n}D_j$, 而 $\bigcup_nA_n=\bigcup_nD_n$. 可数可加性把左边变为级数之和, 有限部分和正是 $\Prob(A_n)$. 递减情形对补集使用递增结论并用补集公式.

一般并集令 $D_1=A_1$, $D_n=A_n\setminus\bigcup_{j<n}A_j$. 它们互斥且有相同并集, 又 $D_n\subseteq A_n$. 所以 $\Prob(\bigcup_nA_n)=\sum_n\Prob(D_n)\leq\sum_n\Prob(A_n)$.
\end{proof}
可数可加性中的``互斥''和``可数''都有作用. 非互斥事件要扣除重复部分. 不可数并集则不能照搬这条公理: 在 $[0,1]$ 的均匀模型中每个单点概率为零, 全部单点的并集却有概率一. 同样, 可数无限集合上不存在各点具有相同概率的均匀概率: 相同值若为正, 总和无限; 若为零, 总和为零.

\originalsection{等可能模型与计数概率}
在有限空间 $\Omega=\{\omega_1,\ldots,\omega_N\}$ 上, 任意非负数字 $p_i$ 满足 $\sum_i p_i=1$ 都确定概率
\[
 \Prob(A)=\sum_{\omega_i\in A}p_i.
\]
反之, 公理与有限可加性说明每个概率都由各单点的概率唯一确定. 若所有点等可能, $p_i=1/N$, 因而 $\Prob(A)=|A|/|\Omega|$. 均匀性是模型假设, 不能由``列出了所有结果''推出来.

\begin{example}
独立公平的两枚八面骰各标 $1,\ldots,8$. 求点数和为五的概率. 另求六次公平抛硬币恰有四次正面的概率.
\end{example}
\begin{solution}
骰子的有序结果有 $8^2=64$ 个, 和为五的有 $(1,4),(2,3),(3,2),(4,1)$ 四个, 所以概率为 $1/16$. 硬币序列有 $2^6$ 个, 正面位置的选择有 $\binom64$ 个, 所以概率为 $\binom64/2^6=15/64$. 这里用到的独立公平假设就是各完整序列等概率; 下一章会把独立性正式定义.
\end{solution}

\begin{example}
有 $m$ 个有标号的记录, 每个记录独立且均匀地取 $d$ 个可能日期之一. 求至少两个记录日期相同的概率. 当 $m\leq d$ 时, 给出易算的近似.
\end{example}
\begin{solution}
至少一次重复的直接分类很多, 所以先数无重复. 所有分配有 $d^m$ 个, 无重复分配有 $d(d-1)\cdots(d-m+1)$ 个, 因而
\[
 \Prob(\text{重复})=1-\prod_{j=0}^{m-1}\left(1-\frac jd\right).
\]
当 $m>d$ 时必有重复, 概率为一. 当各 $j/d$ 很小时, 对数展开 $\log(1-x)=-x+O(x^2)$ 给
\[
 \log\Prob(\text{无重复})
 =-\frac{m(m-1)}{2d}+O\left(\frac{m^3}{d^2}\right).
\]
因此在误差项小时, 无重复概率近似为 $\exp(-m(m-1)/(2d))$. 这解释了为什么发生碰撞所需的记录数量在 $\sqrt d$ 的尺度上, 而不在 $d$ 的尺度上. 用生日解释模型时, 独立性、各日期均匀和忽略闰日都只是简化假设.
\end{solution}

\begin{example}
每条记录独立且均匀地取得 $q$ 个代码之一. $q$ 条记录都不取得某个指定代码的概率是多少? 当 $q$ 增大时有何极限?
\end{example}
\begin{solution}
每个位置有 $q-1$ 个允许代码, 故合格序列有 $(q-1)^q$ 个, 总序列有 $q^q$ 个. 概率为 $(1-1/q)^q$. 因 $q\log(1-1/q)\to-1$, 其极限为 $\ee^{-1}$. 这和下一节错排极限相同, 但两种模型不同: 此处各记录允许重复代码, 错排则始终是双射.
\end{solution}

\originalsection{容斥原理与错排}
对三个事件, 加上单个事件概率, 减去两两交集, 再加回三重交集, 才能让每个结果恰好被计算一次. 一般公式保留了同样的结构.

\begin{theorem}[容斥原理]
任意有限事件族 $A_1,\ldots,A_n$ 满足
\[
 \Prob\left(\bigcup_{i=1}^nA_i\right)
 =\sum_{r=1}^n(-1)^{r+1}
 \sum_{1\leq i_1<\cdots<i_r\leq n}
 \Prob(A_{i_1}\cap\cdots\cap A_{i_r}).
\]
若各 $A_i$ 是有限集合, 将概率替换为元素个数也成立.
\end{theorem}
\begin{proof}
把样本空间按每个 $A_i$ 是否发生分成至多 $2^n$ 个互斥区域. 在一个恰好属于 $m\geq1$ 个事件的区域, 右边的单事件和计入它 $\binom m1$ 次, 二重交集和计入它 $\binom m2$ 次, 依次类推. 总系数为
\[
 \sum_{r=1}^m(-1)^{r+1}\binom mr
 =1-(1-1)^m=1.
\]
不属于任何事件的区域被计入零次. 因而右边对并集内每一区域恰好计入一次, 对外部区域不计入. 对这些互斥区域使用有限可加性便得概率公式. 同样逐个元素计数便得集合公式.
\end{proof}

\begin{example}
将 $n$ 个不同物品均匀随机地一一分配到 $n$ 个有标号的位置. 物品 $i$ 的原位置为 $i$. 求没有一个物品回到原位置的概率, 并求恰有 $k$ 个回到原位置的概率.
\end{example}
\begin{solution}
令 $A_i=\{\sigma(i)=i\}$. 指定 $r$ 个位置都固定之后, 剩下 $n-r$ 个位置可任意排列, 因而任意指定的 $r$ 重交集概率为 $(n-r)!/n!$. 容斥和中这样的交集有 $\binom nr$ 个, 其总贡献绝对值为
\[
 \binom nr\frac{(n-r)!}{n!}=\frac1{r!}.
\]
对并集取补, 得无固定点概率与错排数
\[
 \Prob(\text{无固定点})=\sum_{r=0}^n\frac{(-1)^r}{r!},\qquad
 D_n=n!\sum_{r=0}^n\frac{(-1)^r}{r!}.
\]
这里 $D_0=1$, $D_1=0$. 指定恰好哪 $k$ 个位置固定有 $\binom nk$ 种, 其余位置必须错排, 所以
\[
 \Prob(\text{恰有 }k\text{ 个固定点})=\frac{\binom nkD_{n-k}}{n!}.
\]
尤其 $k=n-1$ 不可能: 唯一剩余物品也必须回到原位置. 指数函数的 Taylor 级数给 $\sum_{r\geq0}(-1)^r/r!=\ee^{-1}$; 交错级数余项估计给
\[
 \left|\frac{D_n}{n!}-\ee^{-1}\right|\leq\frac1{(n+1)!}.
\]
故无固定点概率快速趋于 $\ee^{-1}$. 把每个固定事件当独立事件, 得到 $(1-1/n)^n$, 虽有同一极限, 却不是有限 $n$ 的正确公式.
\end{solution}

\originalsection{综合练习与解答}
\begin{exercise}
将十二名不同学生分给三个有名称的实验组, 人数分别为五、四、三. 再将十二人分成三个无名称的四人组. 两种分法各有多少种?
\end{exercise}
\begin{solution}
第一种为 $12!/(5!4!3!)$. 第二种先赋予组名, 得 $12!/(4!)^3$, 每个无名称分组有 $3!$ 种赋名方式, 所以为 $12!/((4!)^3 3!)$. 第一种人数不相同且组已有名称, 无须除以 $3!$.
\end{solution}

\begin{exercise}
四种颜色各有六个不同编号的球, 不计顺序随机取八个. 求颜色数量为 $3,3,1,1$ 的概率, 以及至多出现两种颜色的概率.
\end{exercise}
\begin{solution}
总数为 $\binom{24}{8}$. 第一个事件先选哪两种颜色各出现三次, 有 $\binom42$ 种, 然后在各色内选球, 所以概率为
\[
 \frac{\binom42\binom63^2\binom61^2}{\binom{24}{8}}.
\]
第二个事件不可能只出现一种颜色, 因为每色只有六球. 所以每个合格集合恰好对应唯一的颜色对, 合格数为 $\binom42\binom{12}{8}$, 概率为其除以 $\binom{24}{8}$. 若取球数改为不超过六, 只出现一种颜色的集合就会被多个颜色对重复计入, 必须另作修正.
\end{solution}

\begin{exercise}
求 $a+b+c=10$ 的非负整数解数量, 其中 $a\leq4,b\leq4$. 另说明在一副每个点数有四张、共有十种点数的牌中, 无序抽五张恰好形成一个三张同点数和一个两张同点数的概率.
\end{exercise}
\begin{solution}
不加限制的解数为 $\binom{12}{2}$. 违反 $a\leq4$ 的解令 $a'=a-5$, 得 $\binom72$ 个, $b$ 的违反数相同. 同时违反时令 $a'=a-5,b'=b-5$, 只有 $a'=b'=c=0$ 一个解. 容斥得 $66-21-21+1=25$.

牌的总手数为 $\binom{40}{5}$. 先指定三张组的点数有十种, 再指定两张组的不同点数有九种, 在各自四张牌内选择, 得概率
\[
 \frac{10\cdot9\binom43\binom42}{\binom{40}{5}}.
\]
三张组与两张组的角色不同, 不应再除以二. 同一原则可计算两个对子: 先无序选择两种点数, 再各选两张, 最后从其余八种点数的三十二张牌选单张, 得 $\binom{10}{2}\binom42^2\cdot32/\binom{40}{5}$.
\end{solution}

\begin{exercise}
四种颜色各有六张不同编号的卡片, 无序取四张. 求至多出现两种颜色的手数. 如果另规定一种获胜模式为``同色且编号恰为 $1,2,3,4$'', 这种模式有多少手?
\end{exercise}
\begin{solution}
先对六个颜色对分别数只使用该颜色对的四张集合, 得 $\binom42\binom{12}{4}$. 只用一种颜色的集合, 在包含这种颜色的三个颜色对中各计一次, 应从三次改为一次. 因而须减去 $2\cdot4\binom64$, 答案为
\[
 \binom42\binom{12}{4}-8\binom64.
\]
这类消重与上一道取八球的题不同: 取四张时单色手牌确实存在. 指定编号的获胜模式每色只有一手, 共有四手. 其概率分别是上述手数或四除以 $\binom{24}{4}$.
\end{solution}
